\chapter{Bilangan dan Himpunan Bilangan}\label{ch:g10:numbers}

Matematika bermula dari bilangan, dan tidak semua bilangan sejenis:
bilangan untuk membilang, bilangan negatif, pecahan, serta bilangan seperti
$\sqrt 2$ atau $\pi$ yang tidak dapat dinyatakan oleh pecahan mana pun. Bab
ini menata semuanya menjadi keluarga yang saling bersarang, memperkenalkan
\omterm{def:g10:numbers:interval}{interval} untuk menyatakan bagian dari garis bilangan, dan memakai \omterm{def:g10:numbers:abs}{nilai
mutlak} untuk mengukur jarak antara dua bilangan.

\section{Keluarga-keluarga bilangan}

\begin{definition}[Himpunan bilangan]\label{def:g10:numbers:sets}
\begin{itemize}
  \item $\N$ adalah himpunan \emph{bilangan cacah}\index{bilangan cacah}:
        $0, 1, 2, 3, \dots$
  \item $\Z$ adalah himpunan \emph{bilangan bulat}\index{bilangan bulat}:
        $\dots, -2, -1, 0, 1, 2, \dots$
  \item $\Q$ adalah himpunan \emph{bilangan rasional}\index{bilangan rasional}:
        semua hasil bagi $\frac{p}{q}$ dengan $p \in \Z$, $q \in \N$ dan
        $q \neq 0$.
  \item $\R$ adalah himpunan \emph{bilangan real}\index{bilangan real}: semua
        bilangan yang dapat ditempatkan pada garis bilangan.
\end{itemize}
\end{definition}

\begin{notation}
Lambang $\in$ dibaca ``anggota dari'': $3 \in \N$, $-5 \in \Z$,
$\frac23 \in \Q$. Lambang $\subset$ dibaca ``termuat dalam'': setiap
\omterm{def:g10:numbers:sets}{bilangan cacah} adalah \omterm{def:g10:numbers:sets}{bilangan bulat}, setiap \omterm{def:g10:numbers:sets}{bilangan bulat} adalah \omterm{def:g10:numbers:sets}{bilangan
rasional} (misalnya $-5 = \frac{-5}{1}$), dan setiap \omterm{def:g10:numbers:sets}{bilangan rasional}
adalah \omterm{def:g10:numbers:sets}{bilangan real}, sehingga
\[
\N \subset \Z \subset \Q \subset \R .
\]
Sebuah garis miring meniadakan lambang: $\frac12 \notin \Z$.
\end{notation}

\begin{omfigure}
\begin{tikzpicture}[scale=0.9]
  \draw[omExo, thick] (0,0) ellipse (5.4 and 2.9);
  \draw[omProp, thick] (-0.8,0) ellipse (4.0 and 2.2);
  \draw[omThm, thick] (-1.6,0) ellipse (2.6 and 1.5);
  \draw[omDef, thick] (-2.4,0) ellipse (1.3 and 0.8);
  \node[omDef, font=\small\bfseries] at (-2.4,0.45) {$\N$};
  \node[omThm, font=\small\bfseries] at (-1.6,1.1) {$\Z$};
  \node[omProp, font=\small\bfseries] at (-0.8,1.8) {$\Q$};
  \node[omExo, font=\small\bfseries] at (0,2.5) {$\R$};
  \node[font=\small] at (-2.4,-0.15) {$0 \quad 7$};
  \node[font=\small] at (-0.4,-0.9) {$-5$};
  \node[font=\small] at (1.0,-0.4) {$\tfrac23 \quad {-1.7}$};
  \node[font=\small] at (3.9,-0.6) {$\sqrt2 \quad \pi$};
\end{tikzpicture}

{\small Keluarga bilangan yang saling bersarang: $\N \subset \Z \subset \Q \subset
\R$. Setiap cincin memuat bilangan yang bukan anggota keluarga yang lebih kecil.}
\end{omfigure}

\begin{example}\label{ex:g10:numbers:classify}
Mari kita tempatkan beberapa bilangan pada keluarga terkecil yang memuatnya.
\begin{itemize}
  \item $\frac{28}{4} = 7$, jadi $\frac{28}{4} \in \N$ meskipun ditulis
        sebagai pecahan: sederhanakan lebih dahulu.
  \item $-3.5 = -\frac{35}{10} = -\frac72$ bersifat rasional tetapi bukan \omterm{def:g10:numbers:sets}{bilangan bulat}.
  \item $0.333\dots$ (angka $3$ berulang tanpa henti) sama dengan $\frac13$,
        yaitu sebuah \omterm{def:g10:numbers:sets}{bilangan rasional}.
  \item $\sqrt 2$ dan $\pi$ bersifat real tetapi tidak rasional, seperti
        yang segera kita lihat untuk $\sqrt 2$; bilangan semacam itu disebut
        \emph{irasional}\index{bilangan irasional}.
\end{itemize}
\end{example}

\begin{theorem}[Sifat irasional $\sqrt 2$]\label{thm:g10:numbers:sqrt2}
Bilangan $\sqrt 2$ tidak rasional: tidak ada pecahan \omterm{def:g10:numbers:sets}{bilangan bulat} yang kuadratnya
$2$.
\end{theorem}

\begin{proof}
Kita bernalar dengan kontradiksi, selangkah demi selangkah.
\begin{enumerate}
  \item Andaikan $\sqrt 2 = \frac{p}{q}$ dengan $p$ dan $q$ \omterm{def:g10:numbers:sets}{bilangan bulat}
        positif, dan pecahan itu sudah paling sederhana, sehingga $p$ dan $q$
        tidak keduanya genap.
  \item Mengkuadratkan kedua ruas memberi $2 = \frac{p^2}{q^2}$, yaitu
        $p^2 = 2q^2$. Jadi $p^2$ genap.
  \item Jika $p$ ganjil, katakan $p = 2k+1$, maka
        $p^2 = 4k^2 + 4k + 1$ akan ganjil. Karena $p^2$ genap, $p$ pasti
        genap: $p = 2k$ untuk suatu \omterm{def:g10:numbers:sets}{bilangan bulat} $k$.
  \item Substitusi: $(2k)^2 = 2q^2$, jadi $4k^2 = 2q^2$, jadi
        $q^2 = 2k^2$. Dengan alasan yang sama seperti pada langkah 3, $q$ genap.
  \item Kini $p$ dan $q$ keduanya genap, bertentangan dengan langkah 1.
        Pengandaian tadi mustahil: $\sqrt 2$ \omterm{ex:g10:numbers:classify}{irasional}.
\end{enumerate}
\end{proof}

\begin{remark}
\omterm{def:g10:numbers:sets}{Bilangan rasional} tepat adalah bilangan yang bentuk desimalnya berhenti
(seperti $\frac72 = 3.5$) atau akhirnya mengulang blok yang sama tanpa henti
(seperti $\frac13 = 0.333\dots$ atau $\frac{1}{7} = 0.142857\,142857\dots$).
\omterm{ex:g10:numbers:classify}{Bilangan irasional} seperti $\sqrt2 = 1.41421356\dots$ mempunyai angka desimal
yang tidak pernah membentuk pola berulang. Ciri ini kita terima tanpa bukti
pada tingkat ini.
\end{remark}

\section{Garis bilangan dan interval}

\omterm{def:g10:numbers:sets}{Bilangan real} memenuhi sebuah garis: setelah dipilih titik asal $0$ dan satuan
panjang, setiap \omterm{def:g10:numbers:sets}{bilangan real} berpadanan dengan tepat satu titik.

\begin{definition}[Interval]\label{def:g10:numbers:interval}
Misalkan $a$ dan $b$ \omterm{def:g10:numbers:sets}{bilangan real} dengan $a < b$. Sebuah
\emph{interval}\index{interval} adalah himpunan semua \omterm{def:g10:numbers:sets}{bilangan real} di antara
dua batas. Kurung siku berarti batas itu termasuk, kurung biasa berarti batas
itu tidak termasuk:
\begin{center}
\begin{tabular}{c|c}
notasi & keterangan \\
\hline
$\intcc{a}{b}$ & $a \leq x \leq b$ (kedua ujung termasuk) \\
$\intoo{a}{b}$ & $a < x < b$ (kedua ujung tidak termasuk) \\
$\intco{a}{b}$ & $a \leq x < b$ \\
$\intoc{a}{b}$ & $a < x \leq b$ \\
$\intco{a}{+\infty}$ & $x \geq a$ \\
$\intoo{-\infty}{b}$ & $x < b$ \\
\end{tabular}
\end{center}
Lambang $-\infty$ dan $+\infty$ (``takhingga'') bukan bilangan, melainkan
sekadar cara mengatakan bahwa interval itu berlanjut tanpa henti; kurung di
sebelahnya selalu kurung biasa. Seluruh garis $\R$ adalah interval
$\intoo{-\infty}{+\infty}$.
\end{definition}

\begin{omfigure}
\begin{tikzpicture}[scale=1.0]
  \draw[-Stealth] (-4.4,0) -- (4.6,0) node[below, font=\small] {$x$};
  \foreach \x in {-4,...,4} \draw (\x,-0.07) -- (\x,0.07);
  \foreach \x in {-4,-2,0,2,4}
    \node[below, font=\small] at (\x,-0.1) {$\x$};
  \draw[omDef, very thick] (-3,0.02) -- (1,0.02);
  \fill[omDef] (-3,0.02) circle (2.2pt);
  \draw[omDef, fill=white, thick] (1,0.02) circle (2.2pt);
  \node[omDef, above, font=\small] at (-1,0.15) {$\intco{-3}{1}$};
  \draw[omThm, very thick] (2,0.02) -- (4.5,0.02);
  \draw[omThm, fill=white, thick] (2,0.02) circle (2.2pt);
  \node[omThm, above, font=\small] at (3.4,0.15) {$\intoo{2}{+\infty}$};
\end{tikzpicture}

{\small \omterm{def:g10:numbers:interval}{Interval} pada garis bilangan: titik penuh untuk batas yang termasuk,
titik kosong untuk batas yang tidak termasuk.}
\end{omfigure}

\begin{definition}[Irisan dan gabungan]\label{def:g10:numbers:interunion}
Misalkan $I$ dan $J$ dua himpunan \omterm{def:g10:numbers:sets}{bilangan real}. Himpunan bilangan yang
termasuk dalam keduanya disebut \emph{irisan}\index{irisan} $I \cap J$
(dibaca ``$I$ dan $J$''), sedangkan himpunan bilangan yang termasuk dalam
paling sedikit salah satunya disebut \emph{gabungan}\index{gabungan}
$I \cup J$ (dibaca ``$I$ atau $J$'').
\end{definition}

\begin{example}\label{ex:g10:numbers:interunion}
Ambil $I = \intcc{-1}{3}$ dan $J = \intoo{1}{5}$. Gambarlah keduanya pada
garis yang sama: keduanya bertumpang tindih antara $1$ dan $3$. Karena itu
\[
I \cap J = \intoc{1}{3},
\qquad
I \cup J = \intco{-1}{5}.
\]
Perhatikan kurungnya: $1 \notin J$ sehingga $1 \notin I \cap J$, tetapi
$3 \in I$ dan $3 \in J$, sehingga $3 \in I \cap J$.
\end{example}

\begin{method}[Bekerja dengan interval]\label{met:g10:numbers:intervals}
Untuk mencari \omterm{def:g10:numbers:interunion}{irisan} atau \omterm{def:g10:numbers:interunion}{gabungan} dua \omterm{def:g10:numbers:interval}{interval}:
\begin{enumerate}
  \item gambarlah garis bilangan dan tandai keempat batasnya;
  \item arsirlah \omterm{def:g10:numbers:interval}{interval} pertama di atas garis dan \omterm{def:g10:numbers:interval}{interval} kedua di bawahnya;
  \item \omterm{def:g10:numbers:interunion}{irisannya} adalah tempat kedua arsiran bertumpang tindih, \omterm{def:g10:numbers:interunion}{gabungannya}
        adalah tempat paling sedikit satu arsiran hadir;
  \item tentukan setiap kurung dengan menguji apakah batas itu sendiri
        termasuk dalam kedua himpunan (\omterm{def:g10:numbers:interunion}{irisan}) atau paling sedikit dalam satu
        himpunan (\omterm{def:g10:numbers:interunion}{gabungan}).
\end{enumerate}
\end{method}

\section{Nilai mutlak dan jarak}

\begin{definition}[Nilai mutlak]\label{def:g10:numbers:abs}
\emph{Nilai mutlak}\index{nilai mutlak} sebuah \omterm{def:g10:numbers:sets}{bilangan real} $x$ adalah
\[
\abs{x} =
\begin{cases}
x & \text{jika } x \geq 0, \\
-x & \text{jika } x < 0.
\end{cases}
\]
Sebagai contoh $\abs{7} = 7$ dan $\abs{-4} = 4$. Dalam setiap hal berlaku
$\abs{x} \geq 0$.
\end{definition}

\begin{proposition}[Jarak pada garis]\label{prop:g10:numbers:distance}
Untuk semua \omterm{def:g10:numbers:sets}{bilangan real} $a$ dan $b$, jarak antara titik $a$ dan titik
$b$ pada garis bilangan adalah $\abs{b - a}$. Khususnya, $\abs{x}$ adalah
jarak dari $x$ ke $0$.
\end{proposition}

\begin{proof}
Jika $b \geq a$, jarak dari $a$ ke $b$ adalah $b - a \geq 0$, yang sama
dengan $\abs{b-a}$. Jika $b < a$, jaraknya adalah $a - b = -(b - a) > 0$,
yang lagi-lagi sama dengan $\abs{b - a}$ menurut definisi \omterm{def:g10:numbers:abs}{nilai mutlak}.
\end{proof}

\begin{proposition}[Nilai mutlak dan interval]\label{prop:g10:numbers:absinterval}
Misalkan $a$ \omterm{def:g10:numbers:sets}{bilangan real} dan $r > 0$. Maka
\[
\abs{x - a} \leq r
\quad\text{tepat ketika}\quad
x \in \intcc{a - r}{a + r}.
\]
\end{proposition}

\begin{proof}
$\abs{x-a} \leq r$ menyatakan bahwa jarak dari $x$ ke $a$ paling banyak $r$,
yaitu bahwa $x$ terletak tidak lebih jauh daripada $r$ dari $a$ pada kedua
sisi. Bilangan yang memenuhi syarat ini tepat adalah bilangan antara $a - r$
dan $a + r$, dengan kedua batas termasuk.
\end{proof}

\begin{omfigure}
\begin{tikzpicture}[scale=1.1]
  \draw[-Stealth] (-0.8,0) -- (7,0) node[below, font=\small] {$x$};
  \draw[omDef, very thick] (1.5,0.02) -- (5.5,0.02);
  \fill[omDef] (1.5,0.02) circle (2.2pt);
  \fill[omDef] (5.5,0.02) circle (2.2pt);
  \fill (3.5,0) circle (1.6pt);
  \node[below, font=\small] at (3.5,-0.08) {$a$};
  \node[below, font=\small] at (1.5,-0.08) {$a-r$};
  \node[below, font=\small] at (5.5,-0.08) {$a+r$};
  \draw[Stealth-Stealth, omThm] (3.5,0.35) -- (5.5,0.35)
    node[midway, above, font=\small] {$r$};
  \draw[Stealth-Stealth, omThm] (1.5,0.35) -- (3.5,0.35)
    node[midway, above, font=\small] {$r$};
\end{tikzpicture}

{\small Pertidaksamaan $\abs{x - a} \leq r$ menyatakan \omterm{def:g10:numbers:interval}{interval} bilangan yang
berjarak paling banyak $r$ dari $a$.}
\end{omfigure}

\begin{example}\label{ex:g10:numbers:abs}
Selesaikan $\abs{x - 3} \leq 2$. Penyelesaiannya adalah bilangan yang berjarak
paling banyak $2$ dari $3$: \omterm{def:g10:numbers:interval}{interval} $\intcc{1}{5}$. Sebaliknya, \omterm{def:g10:numbers:interval}{interval}
$\intcc{-1}{7}$ berpusat di $\frac{-1+7}{2} = 3$ dan berjari-jari
$\frac{7-(-1)}{2} = 4$, sehingga dapat dinyatakan oleh $\abs{x - 3} \leq 4$.
\end{example}

\section{Hampiran}

\omterm{ex:g10:numbers:classify}{Bilangan irasional}, bahkan sebagian besar pecahan, tidak dapat ditulis secara
tepat dengan angka desimal yang berhingga banyaknya, sehingga dalam praktik
kita menghampirinya.

\begin{definition}[Hampiran dengan ketelitian tertentu]\label{def:g10:numbers:approx}
Bilangan $d$ disebut \emph{hampiran}\index{hampiran} bagi $x$ dengan
ketelitian $10^{-n}$ apabila $\abs{x - d} \leq 10^{-n}$. Pemenggalan dan
pembulatan bentuk desimal pada angka ke-$n$ sama-sama menghasilkan
hampiran seperti itu.
\end{definition}

\begin{example}
Dari $\pi = 3.14159\,26\dots$: pemenggalan $3.141$ dan pembulatan $3.142$
sama-sama merupakan \omterm{def:g10:numbers:approx}{hampiran} bagi $\pi$ dengan ketelitian $10^{-3}$.
Pembulatan berjarak kurang dari $\frac12 \times 10^{-3}$, sedangkan
pemenggalan hanya menjamin $10^{-3}$. Menulis $3.141 \leq \pi \leq 3.142$
berarti \emph{mengapit} $\pi$ di antara dua batas desimal.
\end{example}

\section{Latihan}

\begin{exercise}[$\star$]\label{exo:g10:numbers:1}
Untuk setiap bilangan, sebutkan yang terkecil di antara himpunan $\N$, $\Z$,
$\Q$, $\R$ yang memuatnya:
\[
\frac{15}{3}, \qquad -7, \qquad \frac{22}{7}, \qquad \sqrt{9}, \qquad
\sqrt{10}, \qquad -2.4, \qquad 0 .
\]
\end{exercise}

\begin{exercise}[$\star$]\label{exo:g10:numbers:2}
Tulislah setiap pernyataan berikut dengan notasi \omterm{def:g10:numbers:interval}{interval}, lalu gambarlah pada
garis bilangan: (a) $-2 \leq x < 5$; (b) $x > 3$; (c) $x \leq -1$;
(d) jarak dari $x$ ke $2$ paling banyak $3$.
\end{exercise}

\begin{exercise}[$\star$]\label{exo:g10:numbers:3}
Hitunglah $I \cap J$ dan $I \cup J$ untuk
\[
\text{(a) } I = \intcc{-3}{2},\ J = \intco{0}{4};
\qquad
\text{(b) } I = \intoo{-\infty}{1},\ J = \intco{-2}{+\infty}.
\]
\end{exercise}

\begin{exercise}[$\star$]\label{exo:g10:numbers:4}
Hitunglah tanpa kalkulator:
\[
\abs{-6}, \qquad \abs{4 - 9}, \qquad \abs{-3 - 5}, \qquad
\abs{\sqrt 2 - 1}, \qquad \abs{1 - \sqrt 2}.
\]
\end{exercise}

\begin{exercise}[$\star$]\label{exo:g10:numbers:5}
Selesaikan persamaan dan pertidaksamaan berikut, lalu berikan himpunan penyelesaiannya:
\[
\abs{x} = 5, \qquad \abs{x - 1} = 3, \qquad \abs{x - 4} \leq 1, \qquad
\abs{x + 2} < 3 .
\]
(Perhatikan bahwa $\abs{x+2} = \abs{x - (-2)}$ adalah jarak ke $-2$.)
\end{exercise}

\begin{exercise}[$\star\star$]\label{exo:g10:numbers:6}
Nyatakan setiap \omterm{def:g10:numbers:interval}{interval} berikut dengan pertidaksamaan berbentuk
$\abs{x - a} \leq r$ atau $\abs{x - a} < r$:
\[
\intcc{2}{8}, \qquad \intoo{-5}{1}, \qquad \intcc{-7}{-3}.
\]
\end{exercise}

\begin{exercise}[$\star\star$]\label{exo:g10:numbers:7}
Tunjukkan bahwa $0.272727\dots$ (blok $27$ berulang tanpa henti) bersifat rasional.
(Petunjuk: namai bilangan itu $x$ lalu hitunglah $100x - x$.)
\end{exercise}

\begin{exercise}[$\star\star$]\label{exo:g10:numbers:8}
Benar atau salah? Berikan alasan berupa argumen atau contoh penyangkal.
\begin{enumerate}
  \item Jumlah dua \omterm{def:g10:numbers:sets}{bilangan bulat} adalah \omterm{def:g10:numbers:sets}{bilangan bulat}.
  \item Hasil bagi dua \omterm{def:g10:numbers:sets}{bilangan bulat} adalah \omterm{def:g10:numbers:sets}{bilangan bulat}.
  \item Jumlah dua \omterm{def:g10:numbers:sets}{bilangan rasional} bersifat rasional.
  \item Jumlah sebuah \omterm{def:g10:numbers:sets}{bilangan rasional} dan sebuah \omterm{ex:g10:numbers:classify}{bilangan irasional}
        bersifat \omterm{ex:g10:numbers:classify}{irasional}.
\end{enumerate}
\end{exercise}

\begin{exercise}[$\star\star$]\label{exo:g10:numbers:9}
Dengan memakai $1.414 \leq \sqrt 2 \leq 1.415$, apitlah bilangan $2\sqrt2$,
$\sqrt2 + 3$ dan $-\sqrt2$ di antara dua batas desimal.
\end{exercise}

\begin{exercise}[$\star\star\star$]\label{exo:g10:numbers:10}
Sesuaikan bukti \cref{thm:g10:numbers:sqrt2} untuk menunjukkan bahwa $\sqrt 3$
\omterm{ex:g10:numbers:classify}{irasional}. (Gantilah ``genap'' dengan ``kelipatan $3$'': periksa lebih dahulu
bahwa jika $p^2$ kelipatan $3$, maka $p$ pun demikian, dengan meninjau sisa
$0$, $1$, $2$ dari $p$ pada pembagian oleh $3$.)
\end{exercise}

\section{Soal: Di antara dua bilangan mana pun}

\begin{problem}[{Soal akhir pekan --- bilangan rasional dan irasional
saling berselang-seling: setiap interval, sekecil apa pun, memuat tak hingga
banyak bilangan dari kedua jenis, dan tidak ada pengukuran yang mampu
membedakannya}]
\label{pb:g10:numbers:1}
\omterm{def:g10:numbers:sets}{Bilangan rasional} tampak seperti kerumunan (semua pecahan!) dan \omterm{ex:g10:numbers:classify}{bilangan
irasional} seperti perkecualian yang langka ($\sqrt2$, $\pi$). Soal ini
menyingkap gambaran yang sebenarnya: kedua keluarga itu
\emph{berselang-seling} begitu rapat sehingga setiap \omterm{def:g10:numbers:interval}{interval} garis bilangan,
sekecil apa pun, memuat tak hingga banyak bilangan dari kedua jenis --- dengan
akibat yang mengejutkan, misalnya: tidak ada pengukuran fisis, seberapa pun
telitinya, yang dapat memutuskan apakah suatu panjang bersifat rasional.

\medskip
\noindent\textbf{Bagian I --- Empat kerajaan.}
\begin{enumerate}
  \item Untuk setiap bilangan, sebutkan yang terkecil di antara himpunan
        $\N$, $\Z$, $\Q$, $\R$ yang memuatnya
        (\cref{def:g10:numbers:sets}): $-7$; \ $\frac{13}{4}$; \
        $\sqrt{16}$; \ $0.121212\ldots$ (ingat kembali soal akhir pekan
        tentang desimal berulang pada jilid sebelumnya); \ $\sqrt8$; \
        $\pi$ (terima saja sifat \omterm{ex:g10:numbers:classify}{irasionalnya} --- baru dibuktikan pada
        jilid universitas).
  \item Buktikan bahwa $\Q$ tertutup terhadap penjumlahan dan
        perkalian: jika $x = \frac pq$ dan $y = \frac rs$ rasional,
        tulislah $x + y$ dan $xy$ masing-masing sebagai satu pecahan.
  \item Simpulkan dengan kontradiksi: (a) jumlah sebuah \omterm{def:g10:numbers:sets}{bilangan rasional} dan
        sebuah \omterm{ex:g10:numbers:classify}{bilangan irasional} bersifat \omterm{ex:g10:numbers:classify}{irasional}; (b) hasil kali sebuah
        \omterm{def:g10:numbers:sets}{bilangan rasional} \emph{taknol} dan sebuah \omterm{ex:g10:numbers:classify}{bilangan irasional} bersifat \omterm{ex:g10:numbers:classify}{irasional}.
  \item Tunjukkan bahwa himpunan \omterm{ex:g10:numbers:classify}{bilangan irasional} tidak tertutup terhadap
        \emph{kedua} operasi itu: berikan dua \omterm{ex:g10:numbers:classify}{bilangan irasional} yang
        jumlahnya rasional, dan dua yang hasil kalinya rasional.
  \item Letakkan $\sqrt2 + \sqrt8$ dan $\sqrt2 \times \sqrt8$ di antara
        keempat kerajaan tadi (sederhanakan $\sqrt8$ lebih dahulu,
        dengan metode jilid sebelumnya).
\end{enumerate}

\medskip
\noindent\textbf{Bagian II --- \omterm{def:g10:numbers:sets}{Bilangan rasional} itu padat.}
\begin{enumerate}[resume]
  \item Carilah sebuah \omterm{def:g10:numbers:sets}{bilangan rasional} yang terletak tegas di antara
        $3.47$ dan $3.48$; lalu carilah yang kedua; kemudian jelaskan cara
        menghasilkan sebanyak yang diinginkan (perbesaran pada soal akhir
        pekan tentang tiadanya bilangan berikutnya di jilid sebelumnya, kini
        dengan bukti yang sudah terlihat).
  \item Teorema umumnya. Misalkan $a < b$ dua \omterm{def:g10:numbers:sets}{bilangan real} sebarang, dengan
        selisih $g = b - a > 0$. Pilih $n$ dengan
        $10^{-n} < g$, lalu tinjau kelipatan $10^{-n}$
        (kisi desimal berlangkah $10^{-n}$). Jelaskan mengapa paling
        sedikit satu titik kisi jatuh tegas di antara $a$ dan $b$,
        lalu simpulkan: \emph{setiap \omterm{def:g10:numbers:interval}{interval} yang panjangnya positif
        memuat sebuah \omterm{def:g10:numbers:sets}{bilangan rasional}}.
  \item Tingkatkan kesimpulan itu: setiap \omterm{def:g10:numbers:interval}{interval} semacam itu memuat
        \emph{tak hingga banyak} \omterm{def:g10:numbers:sets}{bilangan rasional}. (Terapkan pertanyaan 7
        sekali lagi, di dalam \omterm{def:g10:numbers:interval}{interval} yang lebih kecil.)
  \item Kini \omterm{ex:g10:numbers:classify}{bilangan irasional}: diberikan $a < b$, ambil \omterm{def:g10:numbers:sets}{bilangan rasional}
        $r$ yang terletak tegas di dalamnya (pertanyaan 7) lalu tinjau
        bilangan $r + \frac{\sqrt2}{10^k}$. Dengan pertanyaan 3, tunjukkan
        bahwa bilangan itu \omterm{ex:g10:numbers:classify}{irasional}, dan bahwa untuk $k$ cukup besar
        bilangan itu masih terletak di dalam \omterm{def:g10:numbers:interval}{interval}: \emph{setiap \omterm{def:g10:numbers:interval}{interval}
        juga memuat tak hingga banyak \omterm{ex:g10:numbers:classify}{bilangan irasional}}.
  \item Dua teka-teki klasik terpecahkan: adakah \omterm{def:g10:numbers:sets}{bilangan real} positif
        terkecil? Adakah \omterm{def:g10:numbers:sets}{bilangan real} ``tepat setelah
        $3$''? Jawablah keduanya dengan teorema pertanyaan 7, dan
        berilah hormat kepada versi masa kecilnya (soal akhir pekan tentang
        tiadanya bilangan berikutnya di jilid sebelumnya).
\end{enumerate}

\medskip
\noindent\textbf{Bagian III --- \omterm{def:g10:numbers:abs}{Nilai mutlak}, geometri dari
$\R$.}
\begin{enumerate}[resume]
  \item Selesaikan, lalu nyatakan penyelesaiannya sebagai \omterm{def:g10:numbers:interval}{interval} atau
        \omterm{def:g10:numbers:interunion}{gabungan} (\cref{prop:g10:numbers:absinterval}):
        $\abs{x - 5} = 2$; \quad $\abs{x - 5} < 2$; \quad
        $\abs{x + 1} \geq 3$.
  \item Sebuah piston harus dibuat pada ukuran $80$ mm dengan
        toleransi $0.05$ mm. Tulislah syarat itu dengan
        \omterm{def:g10:numbers:abs}{nilai mutlak}, lalu sebagai \omterm{def:g10:numbers:interval}{interval}. Dua piston terukur
        $79.97$ dan $80.06$ mm: bagaimana putusannya?
  \item Ketaksamaan segitiga pada garis:
        $\abs{a + b} \leq \abs a + \abs b$. Periksalah pada
        $(a, b) = (3, -5)$ dan $(-2, -7)$, buktikan ketika $a$ dan
        $b$ bertanda sama dan ketika keduanya bertanda berlawanan,
        lalu nyatakan tepat kapan kesamaan berlaku.
  \item Selesaikan $\abs{x - 2} = \abs{x + 4}$ dengan membacanya sebagai
        kesamaan jarak pada garis. Titik mana pada soal akhir pekan tentang
        dua cermin di jilid sebelumnya yang baru saja dihitung di sini, satu
        dimensi lebih rendah?
  \item Sederhanakan $\sqrt{x^2}$ --- dengan cermat. Ujilah rumus itu
        pada $x = 3$ dan $x = -3$, lalu pakailah untuk menyelesaikan
        $x^2 < 9$ dengan \omterm{def:g10:numbers:abs}{nilai mutlak}.
\end{enumerate}

\medskip
\noindent\textbf{Bagian IV --- Apa yang tak dapat diputuskan oleh pengukuran.}
\begin{enumerate}[resume]
  \item Seorang fisikawan mengukur sebatang logam: $1.41421 \pm 0.00001$ m.
        Dapatkah pengukuran semacam itu --- yang ini atau yang lebih teliti
        di kemudian hari --- \emph{membuktikan} bahwa panjang batang itu
        \omterm{ex:g10:numbers:classify}{irasional}? Atau bahwa panjangnya rasional? Pakailah pertanyaan 7
        dan 9 pada \omterm{def:g10:numbers:interval}{interval} toleransinya, lalu jelaskan mengapa hanya
        \emph{bukti} (seperti untuk diagonal, pada soal akhir pekan tentang
        sifat \omterm{ex:g10:numbers:classify}{irasional} di jilid sebelumnya) yang dapat memutuskannya.
  \item Pemenggalan $1.4$, $1.41$, $1.414$, $1.4142,\dots$
        dari $\sqrt2$ semuanya rasional. Apa yang ditunjukkannya tentang
        seberapa rapat $\Q$ mendesak setiap \omterm{ex:g10:numbers:classify}{bilangan irasional}?
        Rumuskan fakta umumnya.
  \item Aritmetika \omterm{def:g10:numbers:interval}{interval}: nilai terukur $a = 2.5 \pm 0.1$
        dan $b = 1.2 \pm 0.1$. Apitlah $a + b$ dan $a \times b$
        di antara batas yang pasti. Operasi mana yang lebih banyak
        menurunkan ketelitian?
  \item Aturan seorang insinyur mengatakan bahwa galat pengukuran yang
        saling bebas ``bertambah'' pada penjumlahan: nyatakan dengan
        ketaksamaan segitiga (pertanyaan 13) mengapa galat $a + b$ paling
        banyak sama dengan jumlah kedua galat --- ketaksamaan itulah aturan
        keinsinyuran tersebut.
  \item Penutup, dalam satu paragraf singkat: rangkailah gambaran garis
        \omterm{def:g10:numbers:sets}{bilangan real} yang dibangun oleh soal ini dan para pendahulunya ---
        tiadanya bilangan berikutnya (soal akhir pekan tentang tiadanya
        bilangan berikutnya di jilid sebelumnya), \omterm{def:g10:numbers:sets}{bilangan rasional} $=$
        desimal berulang (soal akhir pekan tentang desimal berulang di jilid
        sebelumnya), kedua keluarga padat (pertanyaan 8 dan 9), pengukuran
        yang selamanya tak menentukan (pertanyaan 16). Akhiri dengan
        pancingan yang belum dapat dibuktikan bab ini: dalam pengertian yang
        tepat, bilangan \emph{\omterm{ex:g10:numbers:classify}{irasional} jauh lebih banyak} daripada \omterm{def:g10:numbers:sets}{bilangan
        rasional} --- jilid universitas menghitung ketakhinggaan itu.

\end{enumerate}
\end{problem}
